\chapter{九推与推免服务系统填报}
\label{ch09}
\ist{九推与系统填报}
\begin{keybox}
这一章讲的是“九推”——每年 9 月下旬到 10 月下旬，在“全国推荐免试攻读研究生（免初试、转段）信息公开管理服务系统”上完成的注册、填志愿、收复试通知、确认待录取的全部动作。它是整本手册里时效性最强的一章：系统每年只在固定窗口开放，时间轴几乎每年都会微调。以最近一届为例，9 月 18 日起注册并查看本人推免资格，9 月 21 日 9:00 开始填报志愿、10 月 18 日截止填报，9 月 22 日 8:00 起确认复试通知，9 月 24 日 14:00 起确认待录取，10 月 20 日系统关闭。本章所有日期均属 2027 届（2026 年 9 月），换成别的年份就会全部失效；凡引用更早年份的地方，都已明确标注年份。请务必在执行前，用研招网与省级教育招生考试机构的当年公告复核一遍时间与规则，不要拿往年的经验直接操作。
\end{keybox}

\section{复试通知}
\begin{qa}{E0014}{被锁了 48 小时还没收到复试通知，解锁之后还能再填同一所学校吗？}
系统层面允许重填同一所学校，但代价是再锁 48 小时。学信网帮助中心与研招网口径一致：志愿提交 48 小时后仍未接到复试通知、复试未通过，或招生单位在 48 小时内明确拒绝，该志愿自动解锁，可继续填报其他志愿（\nref{S0236}、\nref{S0444}、\nref{S0389}）。规则里没有“同一招生单位只能报一次”的限制。真正的门槛在招生单位一侧：上财、中科大等“按批次先申请先审核、额满为止”（\nref{S0183}）；东华大学会在官网实时更新各学院是否还招（\nref{S0419}）；武大计算机学院更要求预报名考生在全国系统开放后 5 小时内填报（\nref{S0119}）。所以解锁后要不要重填，先看目标院系官网是否仍开放，再看你手里还剩多少时间窗口。
\nref{S0236}
\nref{S0444}
\nref{S0389}
\nref{S0183}
\nref{S0419}
\nref{S0119}
\end{qa}

\iin{教育部学生服务与素质发展中心（学信网/研招网）}
\ith{志愿解锁}
\ist{九推与系统填报}
\begin{qa}{E0015}{三个志愿全填同一所学校的不同学院，可以吗？会不会互相冲突？}
可以，但有校内互斥风险。推免系统层面只约束“同时最多 3 个平行志愿、单个志愿提交后 48 小时内不可改”（\nref{S0444}、\nref{S0389}），同一招生单位的不同院系各占一个志愿并不违规；清华经管就明确“报名两个院系将占用两个志愿”（\nref{S0254}）。风险来自校内规则：不少学院限报一个志愿，如人大财政金融学院“同时申请两个以上专业者，材料不予受理”（\nref{S0019}）、北大智能学院“申请者只能申请院内一个志愿”（\nref{S0299}）、华东师大信电学院“只接受每位学员申请我院其中一个专业/方向”（\nref{S0217}）。把三个志愿全押在一所学校，等于放弃 48 小时锁定期内的全部横向机会。建议同校最多占 1–2 个志愿，其余留给不同单位。
\nref{S0444}
\nref{S0389}
\nref{S0254}
\nref{S0019}
\nref{S0299}
\nref{S0217}
\end{qa}

\iin{教育部学生服务与素质发展中心（学信网/研招网）}
\ith{志愿搭配}
\ist{九推与系统填报}
\begin{qa}{E0016}{志愿提交后想改，必须先撤回吗？撤回有没有次数限制？}
系统里没有“撤回”按钮，规则是硬性的 48 小时内不可修改，因此不存在“撤回次数限制”这个说法。官方口径：推免生（有特殊政策要求的专项计划推免生除外）志愿填报确认后 48 小时内不可以修改；若招生单位在 48 小时内通过系统拒绝推免生申请，或 48 小时后仍未接到复试通知、复试未通过，则该志愿自动解锁，推免生可继续填报其他志愿（\nref{S0236}、\nref{S0444}、\nref{S0389}）。唯一的例外是立功表彰退役军人免试生：可填写三个志愿，报名期间可修改（\nref{S0444}）。所以提交前必须核对招生单位、院系所、专业与研究方向；一旦提交，48 小时内唯一的出路是等对方拒绝或等系统自动解锁。
\nref{S0236}
\nref{S0444}
\nref{S0389}
\end{qa}

\iin{教育部学生服务与素质发展中心（学信网/研招网）}
\ith{志愿撤回}
\ist{九推与系统填报}
\begin{qa}{E0017}{已经收到复试通知了，还能撤回志愿吗？撤了会不会连复试资格一起没？}
不能主动撤回，但收到复试通知不等于志愿解锁。志愿的锁定与解锁只看三件事：48 小时未接通知、复试未通过、或你拒绝招生单位发送的复试/待录取通知——满足任一即自动解锁（\nref{S0236}、\nref{S0182}）。所以只要这个志愿还挂着，你就不能再改填别的；想腾出志愿位，只能拒绝该单位的复试通知或待录取通知。而拒绝是不可逆操作：武大计算机学院明确“考生查看复试通知，并对复试通知进行接受或拒绝操作，一旦操作不可更改”（\nref{S0119}）。还要注意，未在系统内确认接受复试通知的考生不能被招生单位录取（\nref{S0389}），“放着不点”同样会让资格落空。
\nref{S0236}
\nref{S0182}
\nref{S0119}
\nref{S0389}
\end{qa}

\iin{教育部学生服务与素质发展中心（学信网/研招网）}
\ith{志愿撤回}
\ist{九推与系统填报}
\begin{qa}{E0018}{接受了一所学校的复试通知之后，其他志愿还能继续走吗？}
能。接受某校的复试通知只是把“可以参加该校复试”这件事定下来，不锁死其他人。系统真正收紧的时刻是“接受待录取通知”：四川省公告写明“推免生接受招生单位发放的待录取通知后，不能再填报其他志愿，或接受其他招生单位复试、待录取通知，其他招生单位也不能再向该推免生发放复试或待录取通知”（\nref{S0389}）。东华大学的流程说明也强调“若考生未在系统中确认接受复试通知，学院不得向其发送待录取通知”（\nref{S0419}）。所以正确姿势是：复试通知照接，三个志愿继续并行推进，把“接受待录取”这一下留到全部结果出来之后。
\nref{S0389}
\nref{S0419}
\end{qa}

\iin{www.sceea.cn}
\ith{复试通知}
\ist{九推与系统填报}
\begin{qa}{E0019}{复试通知要求在多少小时内确认？超时不确认会有什么后果？}
没有全国统一的确认时限，这一条属于院校特定。国家层面的表述是“应在招生单位规定时间内通过系统回复是否同意参加复试，否则招生单位可取消复试通知”（\nref{S0389}、\nref{S0444}），具体小时数由招生单位在通知里写死。各校差异很大：大连理工商学院要求“在我院发出复试通知后 4 小时内确认”（\nref{S0550}、\nref{S0106}）；西南财经大学、中南大学直博为 24 小时（\nref{S0483}、\nref{S0625}）；中国社科院大学经济学院为 2 小时（\nref{S0308}）。后果只有两条：超时未回复，招生单位可取消复试通知；且未在系统中确认接受复试通知的考生不能被招生单位录取（\nref{S0389}）。收到通知当刻就确认，不要攒着。
\nref{S0389}
\nref{S0444}
\nref{S0550}
\nref{S0106}
\nref{S0483}
\nref{S0625}
\nref{S0308}

以该条通知实际送达时间起算，不要按收到短信时间或自己所在时区猜测系统截止点；马上登录系统查看通知详情、具体时限和按钮状态，并设置早于截止时间的提醒。若多条通知冲突，立即联系发送通知的院系研招联系人，确认是否需先确认、能否延时；电话后补发邮件留痕。系统允许确认的年度窗口，不会自动延长学校给个人设置的回复期限。
\end{qa}

\iin{www.sceea.cn}
\ith{复试通知}
\ist{九推与系统填报}
\begin{qa}{E0020}{复试通知是发短信还是只在系统里看？我手机没收到提醒怎么办？}
以系统为准，短信与电话只是辅助。研招网流程明确“推免生提交报考志愿后，招生单位将通过推免服务系统反馈是否同意其参加复试，推免生应及时登录系统，查看相关通知信息”（\nref{S0444}、\nref{S0419}）。部分学校会叠加电话或邮件，如哈工大“将通过全国推免服务系统、电话或电子邮件等方式向考生发送复试相关信息”（\nref{S0067}），西南财经大学递补阶段“以电话或短信方式通知”（\nref{S0483}），但这些都是各家做法，不构成保证。所以要把系统当成唯一信息源，那几天每天至少查两次，同时保持手机畅通、邮箱可用（\nref{S0577}）。系统内的疑问可直接联系招生单位或省级机构（\nref{S0464}）。
\nref{S0444}
\nref{S0419}
\nref{S0067}
\nref{S0483}
\nref{S0577}
\nref{S0464}
\end{qa}

\iin{教育部学生服务与素质发展中心（学信网/研招网）}
\ith{复试通知}
\ist{九推与系统填报}
\begin{qa}{E0051}{学校官网发的通知和系统里的通知不一致，以哪个为准？}
分两类看。涉及“录取是否生效”的，以系统为准：推免服务系统是全国统一的推免工作信息备案公开平台和网上报考录取系统，资格审核确认、报考、录取及备案公开均须通过该系统进行，“最终推免生名单、录取名单和推免生学籍注册均以推免服务系统备案信息为准”（\nref{S0464}），未经系统确认接受的待录取无效（\nref{S0389}）。涉及校内环节（材料、复试安排、考核时间）的，以招生单位官网通知为准——研招网早年的常见问题就答过：网站提示时间按教育部招考公告发布，学校的具体时间安排请与其联系咨询（\nref{S0523}）。两者冲突时，按官网准备材料，但务必在系统里完成确认动作。
\nref{S0464}
\nref{S0389}
\nref{S0523}
\end{qa}

\iin{教育部学生服务与素质发展中心（学信网/研招网）}
\ith{通知效力}
\ist{九推与系统填报}
\begin{qa}{E0052}{学校在系统外让我签纸质承诺书，这个东西有约束力吗？}
有约束力，但不能替代系统操作。系统外的承诺书、拟录取承诺书是学校与学生之间的约定，内容一般包括按期毕业、不放弃录取、诚信应考等，违反的后果是学校可取消录取资格或不予录取。例如南昌大学基础医学院要求“拟录取的同学需签订《南昌大学基础医学院拟录取研究生承诺书》”（\nref{S0001}）。但录取是否成立只认系统：“推免生未通过系统确认接受‘待录取’通知，则录取无效”（\nref{S0389}）；河海大学也说明“考生在系统内确认接受待录取后，志愿将被锁定，原则上不得变更”（\nref{S0178}）。签之前先读违约条款，别把系统外的承诺当成锁住名额的凭据。
\nref{S0001}
\nref{S0389}
\nref{S0178}
\end{qa}

\iin{南昌大学江西医学院基础医学院}
\ith{纸质材料}
\ist{九推与系统填报}
\begin{qa}{E0053}{什么时候能在系统里查到自己最终的录取状态？}
看四个时间点。一是系统关服：2027 届 10 月 18 日关闭注册、支付、报名功能，10 月 20 日关闭复试及待录取通知的发送与确认功能（\nref{S0604}、\nref{S0464}），此后系统内的状态即最终备案状态。二是拟录取公示：东华大学称 2027 年推免生拟录取名单预计 10 月中旬公示（\nref{S0419}）；中山大学要求 10 月下旬前完成拟录取名单网上公示（\nref{S0167}）；高校公示期一般不少于 7 日（\nref{S0347}、\nref{S0520}）。三是备案：最终名单以推免服务系统备案信息为准（\nref{S0464}）。四是录取通知书，与系统状态无关，一般与统考生一并发放（\nref{S0178}）。
\nref{S0604}
\nref{S0464}
\nref{S0419}
\nref{S0167}
\nref{S0347}
\nref{S0520}
\nref{S0178}
\end{qa}

\iin{河北大学}
\ith{录取查询}
\ist{九推与系统填报}
\begin{qa}{E0058}{系统里手滑点了“拒绝”，还能恢复吗？}
基本不能恢复，这是本章最不该试错的一步。招生单位侧口径高度一致：武大计算机学院“对复试通知进行接受或拒绝操作，一旦操作不可更改”（\nref{S0119}）；武汉理工“考生接受待录取通知后，原则上不得更改”（\nref{S0577}）；河海大学“确认接受待录取后，志愿将被锁定，原则上不得变更”（\nref{S0178}）。拒绝的直接后果是该志愿自动解锁、可继续填报其他志愿（\nref{S0182}），但该单位通常不会再向你发放通知。若确实误操作，唯一路径是立刻电话联系该招生单位研招办说明情况并请求重新发放通知——能否成功完全取决于对方是否还有名额，系统层面没有撤销入口。
\nref{S0119}
\nref{S0577}
\nref{S0178}
\nref{S0182}
\end{qa}

\iin{武汉大学}
\ith{误操作}
\ist{九推与系统填报}
\begin{qa}{E0060}{是不是在系统里接受了待录取就等于定下来了？后面还会变吗？}
对“择校”来说基本定死，但不是绝对保险。定死的部分：接受待录取后不能再填报其他志愿、接受其他单位的复试或待录取通知（\nref{S0389}）；中山大学明确“在推免服务系统上完成待录取确认的推免生不得变更录取志愿，不得撤销待录取”（\nref{S0167}），华南理工同样表述（\nref{S0479}）。仍可变的部分：录取资格可能因其他原因被取消，常见情形包括入学前未取得本科毕业证书或学士学位、思想政治素质和品德考核不合格、弄虚作假或违纪、被本科所在学校取消推免资格（\nref{S0520}、\nref{S0178}）。所以接受之后要盯的是按期毕业、成绩与纪律不出问题。
\nref{S0389}
\nref{S0167}
\nref{S0479}
\nref{S0520}
\nref{S0178}
\end{qa}

\iin{www.sceea.cn}
\ith{待录取效力}
\ist{九推与系统填报}
\section{待录取}
\begin{qa}{E0010}{推免的平行志愿跟高考平行志愿是一回事吗？会按顺序录取吗？}
不是一回事。高考平行志愿的核心是“分数优先、遵循志愿、一轮投档”，系统按你填的顺序检索并投档，投出去即结束。推免系统是“三个平行志愿（不分主次）”，官方表述为“不区分先后顺序”（\nref{S0444}、\nref{S0389}），不存在按顺序检索投档的机制：三个志愿同时挂在系统里，各招生单位都能下载你的报名信息并独立决定是否发复试通知（\nref{S0182}、\nref{S0236}），可能同时给你发通知；最终由你选择接受哪一个待录取，且只能接受一个（\nref{S0419}）。所以填推免志愿不需要设“冲稳保”梯次来规避投档失败，而应按“想去程度 × 拿通知概率”来配比。
\nref{S0444}
\nref{S0389}
\nref{S0182}
\nref{S0236}
\nref{S0419}
\end{qa}

\iin{教育部学生服务与素质发展中心（学信网/研招网）}
\ith{平行志愿}
\ist{九推与系统填报}
\begin{qa}{E0011}{三个志愿是同时生效还是先填的先处理？先填的学校会先看到我吗？}
系统层面三个志愿同时生效；处理快慢由招生单位掌握。提交后三个志愿都进入已提交状态，招生单位可下载报名信息并决定是否允许参加复试（\nref{S0444}、\nref{S0182}），这里没有“先填先排队”的系统机制。但不少院校自设先到先得：上财“按批次先申请先审核的原则，择优选拔，额满为止”（\nref{S0183}）；大连理工“本着先申请先审核的原则，择优选拔，额满为止”（\nref{S0550}）；东华大学在官网实时更新各学院“继续招收/推免名额已满”（\nref{S0419}）。进入待录取阶段后，武汉理工甚至“按照考生确认待录取通知的时间顺序确定接收名单，额满为止”（\nref{S0577}）。结论：别指望排队，但要早填、早确认。
\nref{S0444}
\nref{S0182}
\nref{S0183}
\nref{S0550}
\nref{S0419}
\nref{S0577}
\end{qa}

\iin{教育部学生服务与素质发展中心（学信网/研招网）}
\ith{平行志愿}
\ist{九推与系统填报}
\begin{qa}{E0021}{待录取通知的确认时限是几小时？半夜发来的那种从什么时候开始算？}
全国没有统一小时数，这一条完全是院校特定；“从何时起算”同样由招生单位在通知里写明。国家层面只说“应在招生单位规定的时间内通过系统答复是否接受录取通知，否则招生单位可取消待录取通知”（\nref{S0389}、\nref{S0444}）。实际差异极大：南京大学商学院要求收到通知后 30 分钟内确认，并规定“9 月 25 日 14:00 前收到通知的以 14:00 为起始时间，14:00 后收到的以收到通知时间为起始时间”（\nref{S0282}）；大连理工部分学院为 1 小时（\nref{S0550}）；中国社科院大学经济学院 2 小时（\nref{S0308}）；西南财经大学、中南大学 24 小时（\nref{S0483}、\nref{S0625}）；华南理工“最长不超过 24 小时”（\nref{S0479}）。
\nref{S0389}
\nref{S0444}
\nref{S0282}
\nref{S0550}
\nref{S0308}
\nref{S0483}
\nref{S0625}
\nref{S0479}
\end{qa}

\iin{www.sceea.cn}
\ith{待录取确认}
\ist{九推与系统填报}
\begin{qa}{E0022}{待录取通知点了接受之后还能反悔吗？系统上有没有撤回入口？}
系统上没有撤回入口，官方口径是“不可更改、不得撤销”。中国药科大学：“考生一旦在系统中确认待录取通知后则不可更改，不可填报其他志愿及接受其他各类通知”（\nref{S0418}）；中山大学：“在推免服务系统上完成待录取确认的推免生不得变更录取志愿，不得撤销待录取”（\nref{S0167}）；华南理工：“与我校达成待录取状态的推免生不得撤销待录取”（\nref{S0479}）；东华大学：“推免生只能接受一个待录取通知，确认接受待录取通知后不能再填报其他志愿”（\nref{S0419}）。河海的表述留了一点口子——“原则上不得变更”（\nref{S0178}），但能否松动取决于学校，不是你能调用的机制。
\nref{S0418}
\nref{S0167}
\nref{S0479}
\nref{S0419}
\nref{S0178}
\end{qa}

\iin{中国药科大学}
\ith{待录取撤回}
\ist{九推与系统填报}
\begin{qa}{E0023}{接受一个待录取之后，其他志愿是不是自动作废？}
比“作废”更强：整条接收通道都会对你关闭。四川省公告原文：“推免生接受招生单位发放的待录取通知后，不能再填报其他志愿，或接受其他招生单位复试、待录取通知，其他招生单位也不能再向该推免生发放复试或待录取通知”（\nref{S0389}）。也就是说，其他志愿不再接通知、你也不能再接受任何新通知，双向封闭。东华大学同样说明，接受待录取后“不能再填报其他志愿，或接受其他复试、待录取通知”（\nref{S0419}）。所以正确顺序是：让多个复试并行推进，等所有想去的单位结果都出来，再决定接受哪一个待录取——接受的那一刻就是结束。
\nref{S0389}
\nref{S0419}
\end{qa}

\iin{www.sceea.cn}
\ith{待录取后果}
\ist{九推与系统填报}
\begin{qa}{E0024}{能不能同时接受两个学校的待录取通知？系统允许吗？}
不允许，系统层面只能接受一个。官方原话：“推免生只能接受一个待录取通知，确认接受待录取通知后不能再填报其他志愿，或接受其他复试、待录取通知”（\nref{S0419}）；中国药科大学表述为“考生只可同意接受一个招生单位一个专业的待录取通知”（\nref{S0418}）。因此不存在“先占两个再挑一个”的操作——一旦接受第一个，第二个待录取通知根本发不进来；如果第二家先发给你，你就面临只能留一个的局面。真正的多选题在“待录取”之前的阶段：三个平行志愿可以同时接收复试通知（\nref{S0444}），把选择权保留到最后。
\nref{S0419}
\nref{S0418}
\nref{S0444}
\end{qa}

\iin{yjszs（高校官网）}
\ith{待录取唯一性}
\ist{九推与系统填报}
\begin{qa}{E0025}{系统里的“待录取”和“拟录取”到底是不是一回事？}
基本是同一状态的两个说法，但含义不同。系统里的按钮只有“待录取”：通过复试的推免生将收到招生单位通过系统发送的“待录取”通知，接受后即成为该招生单位的待录取考生（\nref{S0389}）。“拟录取”是学校口径，指校内已确定、待公示与上级审核的名单，如中南大学“考生确认待录取后，即为我校拟录取直博生”（\nref{S0625}）、华南理工“获得拟录取资格的申请人，必须……接受待录取”（\nref{S0479}）。效力上两者都不是最终录取：最终录取名单和学籍注册均以推免服务系统备案信息为准（\nref{S0464}），入学前还要通过资格复审（\nref{S0520}）。
\nref{S0389}
\nref{S0625}
\nref{S0479}
\nref{S0464}
\nref{S0520}
\end{qa}

\iin{www.sceea.cn}
\ith{待录取与拟录取}
\ist{九推与系统填报}
\begin{qa}{E0026}{9 月 24 日 14:00 才开始确认待录取通知，那之前学校发的算数吗？}
算数——只是你此前点不了“接受”。2027 届的时间轴是错开的：9 月 24 日 8:00 起各招生单位可开始发送待录取通知（\nref{S0347}、\nref{S0489}、\nref{S0067}），9 月 24 日 14:00 起考生才可确认待录取通知（\nref{S0464}、\nref{S0389}）。河北大学的表述最清楚：“2026 年 9 月 24 日 8:00 开通发送待录取通知功能；2026 年 9 月 24 日 14:00 开通推免生确认待录取通知功能”（\nref{S0604}）。所以 8:00 至 14:00 之间收到的通知会先存在系统里，等 14:00 后统一确认；但各校自设的回应时限仍按其通知规定的时间起算（\nref{S0282}），别以为 14:00 前的时间不算。
\nref{S0347}
\nref{S0489}
\nref{S0067}
\nref{S0464}
\nref{S0389}
\nref{S0604}
\nref{S0282}
\end{qa}

\iin{武汉大学}
\ith{待录取时间轴}
\ist{九推与系统填报}
\begin{qa}{E0061}{保研成功之后还能不能报名考研？会不会冲突？}
两者冲突，已被录取就不能再报统考。教育部《2027年全国硕士研究生招生工作管理规定》第二十七条：已被招生单位录取的推免生，不得再报名参加当年全国统一考试、单独考试（\nref{S0442}）；2026 年与 2025 年版规定含义相同（\nref{S0051}、\nref{S0049}）。院校也反复重申，如天津大学“已被我校录取的推免生，不得再报名参加当年全国统一考试、单独考试”（\nref{S0520}），华南理工明确“否则取消推免录取资格”（\nref{S0479}）。反面情形才被允许：“截止规定日期未被招生单位录取的推免生，其推免资格自动失效”（\nref{S0442}），此时可按统考报名时间报考（北京教育考试院口径，\nref{S0352}）。
\nref{S0442}
\nref{S0051}
\nref{S0049}
\nref{S0520}
\nref{S0479}
\nref{S0352}
\end{qa}

\iin{教育部学生服务与素质发展中心（学信网/研招网）}
\ith{推免与统考}
\ist{九推与系统填报}
\section{信息填报}
\begin{qa}{E0003}{个人信息填错了（入学年月、政治面貌这种），已经提交了还能改吗？该找学院还是找系统客服？}
按字段来源分头处理，别只找客服。系统内自动填入的学籍（学历）信息有问题，官方指引是“请尽快向就读高校的学籍管理部门咨询”（\nref{S0389}、\nref{S0444}）——入学年月、学历信息这类字段源头在学籍库，改学籍库才有效。系统操作、账号、支付类问题找研招网/学信网客服或省级机构；对推免招生信息与录取有疑问，研招网说明“考生如有疑问，可直接与相关推荐高校、招生单位联系咨询，或向推荐高校、招生单位所在省级教育行政部门、省级教育招生考试机构申诉”（\nref{S0464}）。志愿提交后 48 小时内不可修改（\nref{S0236}）约束的是志愿，不是个人资料字段。
\nref{S0389}
\nref{S0444}
\nref{S0464}
\nref{S0236}
\end{qa}

\iin{www.sceea.cn}
\ith{信息更正}
\ist{九推与系统填报}
\begin{qa}{E0004}{姓名拼音填错了会影响后面的录取和学籍注册吗？}
影响有限但值得马上改。姓名拼音属于个人资料字段，研招网的填表说明是“按姓名的拼音书写形式填写”（\nref{S0453}）；录取与学籍注册的权威标识是姓名与证件号码，且“最终推免生名单、录取名单和推免生学籍注册均以推免服务系统备案信息为准”（\nref{S0464}）。风险点在于学校会用报名信息制作录取材料，关键字段一旦不符可能被判为无效材料——多个报名系统明确“系统中本人姓名及身份证号码必须准确，提交后不能修改”（\nref{S0114}）。发现拼音错误，先在系统内可编辑项自行更正，改不动就联系就读高校学籍管理部门（\nref{S0389}）。
\nref{S0453}
\nref{S0464}
\nref{S0114}
\nref{S0389}
\end{qa}

\iin{教育部学生服务与素质发展中心（学信网/研招网）}
\ith{信息更正}
\ist{九推与系统填报}
\begin{qa}{E0044}{通讯地址、邮编这些填错了要紧吗？录取通知书按这个寄吗？}
要紧，录取通知书基本按你填的通信地址寄。填表说明把联系方式与档案所在单位地址、邮编并列，档案地址明确“录取后调档使用”（\nref{S0453}）；通知书环节学校也以报名信息为准：中山大学要求考生在指定时间内核对录取通知书邮寄地址，“逾期没有修改邮寄地址的考生，将以报名时所填报的通信地址邮寄”（\nref{S0165}）。统考口径同样强调“通讯地址、联系电话必须详细、准确，且在 2026 年 6 月底前有效”（\nref{S0413}）。录取通知书一般在次年 6 月前后与统考生一并发放（\nref{S0178}）。填一个你能稳定收件的地址，并在学校开放核对窗口时更新一次。
\nref{S0453}
\nref{S0165}
\nref{S0413}
\nref{S0178}
\end{qa}

\iin{教育部学生服务与素质发展中心（学信网/研招网）}
\ith{联系方式}
\ist{九推与系统填报}
\begin{qa}{E0045}{“考生来源”“档案所在地”这些字段到底怎么填？填错影响录取吗？}
按填表说明如实填，档案字段填错确实会影响后续调档。研招网字段说明：“档案所在单位 需如实填写，按国家规定，个人档案不得由个人持有”；“档案所在单位地址”格式为省（区、市）/地市区（盟）/县市旗/乡镇/街村详细地址，需如实填写，最多 80 个字节；“档案所在单位邮政编码”同样“录取后调档使用”（\nref{S0453}）。应届本科生的档案通常就在就读高校，填写前先与学院确认具体到哪一级。录取后学校据此发调档函，填错或填成家庭地址会造成调档延误，进而影响报到与学籍注册（\nref{S0040} 第八条至第十一条）。
\nref{S0453}
\nref{S0040}
\end{qa}

\iin{教育部学生服务与素质发展中心（学信网/研招网）}
\ith{档案字段}
\ist{九推与系统填报}
\begin{qa}{E0046}{政治面貌填错了会不会影响后面的政审？}
直接影响不大，但别把两件事混在一起。系统里的政治面貌“需如实选择”（\nref{S0453}），填错属信息失实，按各校规定可能被要求更正；而录取环节真正的政审是一份单独材料——思想政治素质和品德考核，不合格者不予录取，且会延续到发放录取通知书前的资格复审（\nref{S0178}、\nref{S0520}）。这份材料通常由你档案所在院系党组织填写盖章，例如浙江大学要求政审表“由所在高校相关工作人员认真如实填报”，并加盖二级院系党组织公章方为有效（\nref{S0542}）。要盯的是政审表的填写与盖章，系统里的政治面貌填错请尽早联系学院更正。
\nref{S0453}
\nref{S0178}
\nref{S0520}
\nref{S0542}
\end{qa}

\iin{教育部学生服务与素质发展中心（学信网/研招网）}
\ith{政治面貌}
\ist{九推与系统填报}
\begin{qa}{E0047}{系统里要填成绩排名和绩点吗？填高了学校会核对吗？}
系统个人资料里不强制要求填排名与绩点，但招生单位一定会核。研招网公布的推免生字段清单包含基本信息、户档信息、学习工作经历、发表论文、奖惩、联系方式等类别（\nref{S0453}），成绩排名与绩点并非每项都要求；但校方申请表普遍要求“专业排名证明（需所在学校或院系的教务部门盖章证明）”（\nref{S0514}），北大系里的常见问题明确“总之，以院系盖章数据为准”（\nref{S0264}）。填高或造假的后果写得很直接：各校普遍规定“对在推免过程中弄虚作假……一经查实，将取消其推免资格”（\nref{S0553}），天津大学也把弄虚作假、论文抄袭、虚报获奖列为取消录取资格情形（\nref{S0520}）。
\nref{S0453}
\nref{S0514}
\nref{S0264}
\nref{S0553}
\nref{S0520}
\end{qa}

\iin{教育部学生服务与素质发展中心（学信网/研招网）}
\ith{成绩与排名}
\ist{九推与系统填报}
\begin{qa}{E0048}{退役大学生填系统时“现役军人”那一栏怎么选？军校生该选哪个？}
按研招网填表说明选。官方字段解释：“现役军人 可选择军队在职干部，军校应届本科毕业生（军校学员），非军人，需如实选择”（\nref{S0453}）。所以军校应届本科毕业生选“军校应届本科毕业生（军校学员）”，已在部队任职的选“军队在职干部”，其余选“非军人”。如果你是退役大学生士兵专项计划的考生，还要额外填入伍前所在高校入学年月、批准入伍年月与入伍批准书编号、批准退役年月与退出现役证编号、退役安置地等，专项计划填“退役大学生计划”者这些字段不得为空（\nref{S0453}）；且须在复试前向招生单位提供《入伍批准书》和《退出现役证》（\nref{S0442}）。
\nref{S0453}
\nref{S0442}
\end{qa}

\iin{教育部学生服务与素质发展中心（学信网/研招网）}
\ith{退役与军校}
\ist{九推与系统填报}
\begin{qa}{E0050}{系统客服电话打不通，除了客服还能向谁申诉？}
按问题归属找门路。研招网对推免系统的说明写明了官方渠道：系统内备案公开的推免招生信息均由各推荐高校、招生单位提供，“考生如有疑问，可直接与相关推荐高校、招生单位联系咨询，或向推荐高校、招生单位所在省级教育行政部门、省级教育招生考试机构申诉”（\nref{S0464}）。系统内每个招生单位页面还公示了申诉渠道，包含机构名称、地址、邮编、邮箱、电话、传真与受理部门（\nref{S0465}）。再细分：录取流程与材料问题找招生单位研招办；推荐资格与名额问题找本科学校教务处或学院（\nref{S0553}）；学籍学历信息问题找就读高校学籍管理部门（\nref{S0389}）。
\nref{S0464}
\nref{S0465}
\nref{S0553}
\nref{S0389}
\end{qa}

\iin{教育部学生服务与素质发展中心（学信网/研招网）}
\ith{申诉渠道}
\ist{九推与系统填报}
\section{账号问题}
\begin{qa}{E0001}{推免服务系统什么时候开放注册？我 9 月 18 号进去怎么查不到自己的推免资格？}
开放注册不等于你已具备资格。2027 届的时间轴是：2026 年 9 月 18 日起，推免生通过推免服务系统注册填报个人资料信息并查看本人推免资格；9 月 21 日 9:00 开始填报志愿（\nref{S0464}、\nref{S0389}）。云南师范大学的通知也写明“9 月 18 日系统开通推免生注册功能；9 月 21 日开通推免生报名功能”（\nref{S0568}）。查不到资格，通常是本科学校还没把推荐名单报省级审核并上传系统——名单要先经校内公示（不少于 7 日）与省级审核后才备案（\nref{S0553}、\nref{S0527}）。先问学院教务或推免负责人，别干等系统刷新。
\nref{S0464}
\nref{S0389}
\nref{S0568}
\nref{S0553}
\nref{S0527}
\end{qa}

\iin{教育部学生服务与素质发展中心（学信网/研招网）}
\ith{系统开放}
\ist{九推与系统填报}
\begin{qa}{E0007}{注册时用手机号还是邮箱好？我换号了收不到验证码怎么办？}
注册的是学信网账号，手机号和邮箱都是后续找回凭据，务必都留长期不换的。研招网提醒“请牢记注册手机号、邮箱、证件号码、密码等信息”（\nref{S0523}），统考网报同样要求“考生应牢记实名注册的学信网用户名和密码，后期的信息查询、准考证下载和调剂等环节仍然需要使用”（\nref{S0445}）。换号后收不到验证码属于账号找回问题，公开渠道未见统一的自助流程；可行路径是联系学信网客服，同时告知就读高校学籍管理部门备案（\nref{S0389} 给出了学籍类问题的官方咨询方向）。预防办法只有一个：不要用临时号或校园卡号注册。
\nref{S0523}
\nref{S0445}
\nref{S0389}
\end{qa}

\iin{教育部学生服务与素质发展中心（学信网/研招网）}
\ith{账号注册}
\ist{九推与系统填报}
\begin{qa}{E0008}{学信网账号和推免服务系统的账号是同一个吗？要不要重新注册？}
是同一个，不用重复注册。官方流程第一步就是“注册学信网账号并登录系统：点击注册学信网账号，如已有学信网账号可直接登录系统，登录学信档案后，可核实本人学籍学历信息”（\nref{S0444}）；上海财经大学的说明也写“获得推免资格的同学可注册学信网账号，如已有学信网账号的同学可直接登录系统”（\nref{S0182}）；东华大学提示“推免生需使用本人姓名、证件号码等信息注册学信网账号（已有账号可直接登录）”（\nref{S0419}）。要注意账号必须是你本人实名且已与学籍关联——注册时填错姓名或证件号，会导致登录后看不到本人推免资格。
\nref{S0444}
\nref{S0182}
\nref{S0419}
\end{qa}

\iin{教育部学生服务与素质发展中心（学信网/研招网）}
\ith{账号复用}
\ist{九推与系统填报}
\begin{qa}{E0030}{支付时提示支付失败但银行卡已经扣钱，怎么办？}
先别重复支付，按研招网《2027推免生网上支付报名费注意事项》的“支付查询”处理。原文：如果系统没有提示支付失败或成功，最好通过登录网上银行、电话、ATM、柜台等方式查询账户余额；“如果报名费已经支出，可等待推免服务系统支付状态更新，亦可致电第三方支付平台查询处理”（\nref{S0451}）。只有确认“支付未成功（即银行卡未扣款成功）”时，才需要进入推免服务系统重新支付（\nref{S0451}）。扣款成功而状态未同步的情况，正确动作是等状态刷新并联系第三方支付平台；反复点击支付可能造成重复扣款。
\nref{S0451}
\end{qa}

\iin{教育部学生服务与素质发展中心（学信网/研招网）}
\ith{支付异常}
\ist{九推与系统填报}
\begin{qa}{E0031}{提示银行卡未开通网上支付，除了换卡还有别的支付方式吗？}
有，支付宝可用。研招网《2027推免生网上支付报名费注意事项》列出的方式包括支付宝，以及工行、农行、交行、民生、邮储、中信、华夏、北京银行、兴业、上海银行、招商、建行、光大、浦发、平安等银行，并提示“更多支持的银行请在支付页面跳转的第三方支付平台查询”（\nref{S0451}）。操作要点有三条：一，所持银行卡必须开通网上支付业务功能；二，提前存入报名费，“存入金额一定要大于报名费，而且至少多出 10 元，以免造成支付失败”；三，若支付未成功（银行卡未扣款），需要进入推免服务系统重新支付。
\nref{S0451}
\end{qa}

\iin{教育部学生服务与素质发展中心（学信网/研招网）}
\ith{支付方式}
\ist{九推与系统填报}
\begin{qa}{E0039}{在推免系统里能看到别人的录取公示信息吗？入口在哪？}
能在限定范围内看到。研招网推免系统的说明是：“推免生可在规定时限内查看本人报考并参加复试的招生单位同专业、同层次、同类别推免生录取公示信息”（\nref{S0464}）。注意三个限定：一是须在规定时限内，系统关闭后不再提供；二是范围限于你报考并参加过复试的那家招生单位的同专业、同层次、同类别；三是只含录取公示信息，不含其他考生的报名材料。此外，系统内的推免招生工作办法、推免生名单、拟录取名单等均由各推荐高校和招生单位提供，并已在其官网公示，有疑问可直接联系相关单位或向其所在省级机构申诉（\nref{S0464}）。
\nref{S0464}
\end{qa}

\iin{教育部学生服务与素质发展中心（学信网/研招网）}
\ith{录取公示}
\ist{九推与系统填报}
\begin{qa}{E0043}{推免系统里招生单位的联系方式在哪找？电话打不通还能找谁？}
系统里就能找到官方联系信息。推免服务系统的招生单位页面公示了“申诉渠道”栏目，字段包括机构名称、地址、邮编、邮箱、电话、传真与受理部门（\nref{S0465}），这是最可靠的一手来源。电话打不通时按问题归属换渠道：录取流程、材料审核、复试安排找该招生单位研究生院或研招办；推荐资格、名额、成绩认定找本科学校教务处或学院——北大明确各院系“应畅通学生咨询及申诉渠道，对学生提出的推免相关的意见或申诉应及时调查处理并予以反馈”（\nref{S0553}）；推免政策与系统备案问题可向推荐高校、招生单位所在省级教育行政部门或省级教育招生考试机构申诉（\nref{S0464}）。
\nref{S0465}
\nref{S0553}
\nref{S0464}
\end{qa}

\iin{教育部学生服务与素质发展中心（学信网/研招网）}
\ith{联系方式}
\ist{九推与系统填报}
\begin{qa}{E0049}{登录密码忘了、手机号也换了，怎么找回推免系统账号？}
未见研招网、学信网就“密码遗忘且手机号已更换”给出的公开自助流程，只能走人工渠道。可依据的官方线索有两条：一是账号本身就是学信网账号，研招网要求“考生应牢记实名注册的学信网用户名和密码，后期的信息查询、准考证下载和调剂等环节仍然需要使用”（\nref{S0445}），注册时也提醒“请牢记注册手机号、邮箱、证件号码、密码等信息”（\nref{S0523}）；二是学籍学历类信息问题，官方指引为“尽快向就读高校的学籍管理部门咨询”（\nref{S0389}）。因此建议立刻联系学信网客服走账号申诉，同时告知就读高校教务或学籍部门备案，并务必在填志愿前完成，避免错过 48 小时窗口。
\nref{S0445}
\nref{S0523}
\nref{S0389}
\end{qa}

\iin{教育部学生服务与素质发展中心（学信网/研招网）}
\ith{账号找回}
\ist{九推与系统填报}
\section{志愿填报}
\begin{qa}{E0009}{推免系统一次能填几个志愿？是三个吗？填不满会不会吃亏？}
三个平行志愿，不分主次。研招网操作流程写“可同时填报三个平行志愿（不分主次），每个志愿在提交后的 48 小时内不允许修改”（\nref{S0444}）；四川省 2027 年公告同样表述为“可同时在系统填报 3 个平行志愿（不区分先后顺序）”（\nref{S0389}）；系统首页则说明推免生“可自主多次平行填写多个报考志愿”（\nref{S0464}），意思是三个志愿解锁后可循环使用，并非全程只有三次机会。填不满确实吃亏：少填一个就少一条拿复试通知的入口，而院校普遍“先申请先审核、额满为止”（\nref{S0183}）。注意有特殊政策要求的专项计划推免生受限，如上财规定通过支教团专项计划遴选的推免生只能报考本校专业志愿（\nref{S0182}）。
\nref{S0444}
\nref{S0389}
\nref{S0464}
\nref{S0183}
\nref{S0182}
\end{qa}

\iin{教育部学生服务与素质发展中心（学信网/研招网）}
\ith{志愿数量}
\ist{九推与系统填报}
\begin{qa}{E0027}{10 月 18 日截止填志愿、10 月 20 日结束，中间这两天还能补填或改志愿吗？}
不能补填新志愿，但这两天仍可处理已填志愿产生的通知。研招网列明 2027 届推免生的填报志愿窗口为 9 月 21 日至 10 月 18 日；确认复试通知从 9 月 22 日 8:00 开始，确认待录取通知从 9 月 24 日 14:00 开始，相关流程整体到 10 月 20 日（\nref{S0464}）。因此，10 月 18 日截止的是注册、报名和填报新志愿；18 日至 20 日仍可能收到招生单位发送的复试或待录取通知，也可以按学校要求在系统中确认。河北大学公告明确 10 月 20 日关闭“发送复试考核及待录取通知”以及对应确认功能（\nref{S0604}）。这不是统一延长各校的复试时限，学校自设的确认截止时间仍须遵守。
\nref{S0464}
\nref{S0604}
\end{qa}

\iin{教育部学生服务与素质发展中心（学信网/研招网）}
\ith{填报截止}
\ist{九推与系统填报}
\begin{qa}{E0060}{志愿锁定的 48 小时跨过 10 月 18 日，解锁后还能补填吗？}
不能。2027 届推免服务系统 10 月 18 日截止填报志愿；志愿提交后的 48 小时规则不会把全国填报窗口顺延。即使某个志愿随后因招生单位拒绝或超过锁定时限而解锁，10 月 18 日填报截止后也不能依靠解锁来新增或重填志愿。10 月 18 日至 20 日仍可接收并确认招生单位发送的复试、待录取通知，但系统在 10 月 20 日关闭相关通知的发送与确认功能（\nref{S0464}、\nref{S0604}）。如果当前志愿仍在处理，及时查看系统并直接向招生单位核对进度；不要把“48 小时后可继续填报”的一般规则理解成可以越过当届报名截止日期。
\nref{S0236}
\nref{S0389}
\nref{S0464}
\nref{S0604}
\end{qa}

\iin{教育部学生服务与素质发展中心（学信网/研招网）}
\ith{填报截止}
\ist{九推与系统填报}
\begin{qa}{E0054}{一个志愿能填多个专业方向吗？比如同时报学硕和专硕？}
不能，一个志愿对应一个招生单位的一个院系与专业（方向）。志愿填报要选到招生单位—院系所—专业这一级，部分还要选研究方向甚至导师；清华经管在说明志愿位时写明“报名两个院系将占用两个志愿”（\nref{S0254}），山东大学的校内系统同样是“每个招生单位内仅可选择一个专业报考”（\nref{S0630}）。学硕与专硕的专业代码不同，例如环境科学与工程（083000）是学硕、资源与环境（085700）是专硕，南开大学环境科学与工程学院专门提醒“报名环境工程的同学，务必区分学硕和专硕的代码”（\nref{S0145}），因此同时报学硕与专硕要占两个志愿。
\nref{S0254}
\nref{S0630}
\nref{S0145}
\end{qa}

\iin{清华大学}
\ith{志愿粒度}
\ist{九推与系统填报}
\begin{qa}{E0055}{填志愿时“招生单位”和“院系所”要选到哪一级？选错了能改吗？}
要选到院系所加专业（方向），选错了基本不能当场改。很多学校的接收工作由院系组织、按院系发通知：武大计算机学院要求“通过推免预报名的考生，须于全国推免服务系统开放后的 5 小时内填报我院”（\nref{S0119}）；清华经管的志愿写法是“051经济管理学院-025100（专业学位）金融-01（全日制）金融硕士”（\nref{S0254}），可见院系与专业方向必须与拟录取口径一致。选错的后果是：志愿提交后 48 小时内不可修改（\nref{S0236}），只能等招生单位拒绝或 48 小时后自动解锁再重填（\nref{S0389}）。所以填报前先读该院系的接收办法或实施细则，确认专业与方向写法的准确口径。
\nref{S0119}
\nref{S0254}
\nref{S0236}
\nref{S0389}
\end{qa}

\iin{武汉大学}
\ith{选错院系}
\ist{九推与系统填报}
\begin{qa}{E0056}{专业代码列表里找不到我的专业，选“其他”会不会有问题？}
选“其他”是官方认可的填法，但要按规则补全名称。研招网填表说明明确：“填‘专业代码’时，列表中没有的选择‘其他’；填‘专业名称’时，列表中没有的直接填专业名称，以毕业证书为准”（\nref{S0453}）。北京大学系里的常见问题解答给的是同一口径：“系统里没有我本科专业如何填写？答：可以选择其他，然后填写具体信息”（\nref{S0264}）。风险在于填写的名称若与成绩单、毕业证书不一致可能被判为无效材料，北大明确“系统中所填写的‘所在学校’必须与成绩单上的校名一致”“否则可能被视为无效报名”（\nref{S0264}）。所以选“其他”之后，请按毕业证书上的规范全称填写。
\nref{S0453}
\nref{S0264}
\end{qa}

\iin{教育部学生服务与素质发展中心（学信网/研招网）}
\ith{专业代码}
\ist{九推与系统填报}
\begin{qa}{E0057}{九推填志愿是越早越好吗？早填会不会更早被处理？}
总体上越早越好，但“更早被处理”只在部分学校成立。系统不提供任何“先填先排”的机制，三个志愿同时生效（\nref{S0444}）；时间优势来自院校自设规则：上财、大连理工等“本着先申请先审核的原则，择优选拔，额满为止”（\nref{S0183}、\nref{S0550}）；东华大学在官网实时更新各学院“继续招收”或“推免名额已满”（\nref{S0419}）；进入待录取阶段后，武汉理工明确“学院按照考生确认待录取通知的时间顺序确定接收名单，额满为止”（\nref{S0577}）。所以早填的价值在于抢在名额关闭前进入审核队列，以及早一点拿到通知、多一点比较时间，而不是保证按提交顺序被审。
\nref{S0444}
\nref{S0183}
\nref{S0550}
\nref{S0419}
\nref{S0577}
\end{qa}

\iin{教育部学生服务与素质发展中心（学信网/研招网）}
\ith{填报时机}
\ist{九推与系统填报}
\section{志愿状态}
\begin{qa}{E0002}{系统里“推免资格”那一栏显示未查询到，是不是就等于没保上？还有没有翻盘余地？}
不等于没保上，但你的资格确实还没在系统里生效。研招网说明推免生自 2026 年 9 月 18 日起“注册填报个人资料信息并查看本人推免资格”（\nref{S0464}）；资格以推免服务系统备案信息为准，“未经全国推免服务系统备案的考生不得录取”（\nref{S0347}）。备案的前置条件是本科学校完成校内遴选与公示（不少于 7 日）并报省级机构审核（\nref{S0553}、\nref{S0527}）。所以“未查询到”最常见的原因是学校名单还没上传或省级审核未通过。翻盘方向在校内：立即联系学院推免负责人与教务处核实名单状态，北大也要求各院系畅通咨询与申诉渠道（\nref{S0553}），系统本身无法为你补录资格。
\nref{S0464}
\nref{S0347}
\nref{S0553}
\nref{S0527}
\end{qa}

\iin{教育部学生服务与素质发展中心（学信网/研招网）}
\ith{推免资格}
\ist{九推与系统填报}
\begin{qa}{E0012}{志愿填完是马上被锁定吗？锁定多久才能解锁？}
提交即进入锁定，48 小时内不可修改。学信网帮助中心明确“推免生（有特殊政策要求的专项计划推免生除外）志愿填报确认后 48 小时内不可以修改”（\nref{S0236}）；研招网流程同步表述为“每个志愿在提交后的 48 小时内不允许修改，期间招生单位下载报名信息并决定是否允许考生参加复试”（\nref{S0444}）；四川省 2027 年公告还给出提前解锁的情形：“若招生单位在 48 小时内通过系统拒绝推免生申请，或 48 小时后推免生仍未接到复试通知、复试未通过，则该志愿自动解锁，推免生可继续填报其他志愿”（\nref{S0389}）。也就是说，只要复试通过并拒绝待录取，同样会解锁。
\nref{S0236}
\nref{S0444}
\nref{S0389}
\end{qa}

\iin{教育部学生服务与素质发展中心（学信网/研招网）}
\ith{志愿锁定}
\ist{九推与系统填报}
\begin{qa}{E0013}{锁定 48 小时是系统统一规定，还是各招生单位自己定？}
是系统统一规定，不是各校自定。研招网、学信网帮助中心与省级考试院公告对 48 小时的表述完全一致：志愿提交后的 48 小时内不允许修改；48 小时后仍未接到复试通知、复试未通过，或招生单位在 48 小时内明确拒绝，则该志愿自动解锁（\nref{S0236}、\nref{S0444}、\nref{S0389}、\nref{S0388}）。对比另一套系统可以看清差异：统考调剂服务系统中“每次开放调剂系统持续时间不得低于 12 小时，考生调剂志愿锁定时间不得超过 36 小时”，且明确由招生单位自主设定（\nref{S0049}、\nref{S0457}）。也就是说 36 小时那套是统考调剂规则，不要套用到推免上——推免统一是 48 小时。
\nref{S0236}
\nref{S0444}
\nref{S0389}
\nref{S0388}
\nref{S0049}
\nref{S0457}
\end{qa}

\iin{教育部学生服务与素质发展中心（学信网/研招网）}
\ith{锁定时长}
\ist{九推与系统填报}
\begin{qa}{E0040}{系统显示“该专业已录满”，是不是就不能再填了？}
基本是，名额满了谁也无法录取你。推免招生普遍按“先申请先审核、择优选拔、额满为止”执行（\nref{S0183}、\nref{S0550}），学校还会实时公布各院系的接收情况——东华大学的流程说明就附有“各学院推免生录取情况（实时更新）”，其中已出现“旭日经济与管理学院 推免名额已满”的字样（\nref{S0419}）；未招满的推免名额将用于统招（\nref{S0419}）。所以看到某个专业已满，就直接换目标院系或换学校，别把 48 小时的志愿位浪费在一个已经关闭的专业上。至于系统是否会以“该专业已录满”字样拦截提交，未见官方说明，以你操作时系统的实际提示为准。
\nref{S0183}
\nref{S0550}
\nref{S0419}
\end{qa}

\iin{上海财经大学}
\ith{名额已满}
\ist{九推与系统填报}
\begin{qa}{E0041}{志愿被学校初审拒了，系统会通知我吗？还是只能自己刷？}
系统会变更志愿状态，但不要把“收到通知”当成必然。官方机制是：若招生单位在 48 小时内通过系统拒绝推免生申请，或 48 小时后推免生仍未接到复试通知、复试未通过，则该志愿自动解锁（\nref{S0236}、\nref{S0389}）——拒绝在系统内表现为状态变更与解锁。至于是否另行告知，各校做法不同：哈工大明确会“通过全国推免服务系统、电话或电子邮件等方式”发送相关信息（\nref{S0067}），但这不是普遍承诺；研招网反复强调的是“推免生应及时登录系统，查看相关通知信息”（\nref{S0444}、\nref{S0419}）。所以判断依据是系统状态与解锁，请主动每天查看。
\nref{S0236}
\nref{S0389}
\nref{S0067}
\nref{S0444}
\nref{S0419}
\end{qa}

\iin{教育部学生服务与素质发展中心（学信网/研招网）}
\ith{志愿状态}
\ist{九推与系统填报}
\begin{qa}{E0042}{提交志愿后多久会显示“已查看”？一直停在“已提交”是什么情况？}
官方没有对“已查看”状态给出时限，能确定的是 48 小时这个处理窗口。研招网说明：每个志愿在提交后的 48 小时内不允许修改，“期间招生单位下载报名信息并决定是否允许考生参加复试”（\nref{S0444}、\nref{S0182}）。也就是说招生单位是否下载、何时下载由对方掌握，系统不承诺“多久被查看”。一直停在已提交，到 48 小时会按规则自动解锁（\nref{S0236}），解锁本身就说明该志愿未能产生复试通知。真正该做的不是反复刷新状态，而是在这 48 小时里准备好其他志愿与材料，并尽早填报（\nref{S0419}）。
\nref{S0444}
\nref{S0182}
\nref{S0236}
\nref{S0419}
\end{qa}

\iin{教育部学生服务与素质发展中心（学信网/研招网）}
\ith{志愿状态}
\ist{九推与系统填报}
\section{优营衔接}
\begin{qa}{E0034}{系统里下载的报名信息表、登记表要不要盖章？盖学院章还是教务处章？}
要盖章，但要分清是哪张表、盖哪种章。推免服务系统内生成的是报考与报名信息，真正需要盖章的往往是学校申请表，且各校指定用章不同：北大中文系要求申请的推荐意见“由学校或院系负责人签字并加盖公章”（\nref{S0096}）；南开大学马克思主义学院要求《南开大学2027年推免生报名信息表》“本人签字后扫描”，同时成绩单须“加盖学校教务处公章”、专业排名须“加盖所在学校教务部门公章”（\nref{S0133}）。通用原则是：成绩单与排名证明盖教务部门章，院系申请表盖学院章，推荐信须专家本人签名（\nref{S0264}）。材料不实者不予录取或取消资格（\nref{S0173}）。
\nref{S0096}
\nref{S0133}
\nref{S0264}
\nref{S0173}
\end{qa}

\iin{北京大学}
\ith{材料盖章}
\ist{九推与系统填报}
\begin{qa}{E0035}{有的学校还要在它自己的系统里再报一遍，两个系统都必须填吗？}
是，两个都要走，缺一不可。设有预报名系统的学校普遍规定：申请人在预报名系统报名后，“待教育部开通推免生报名系统后仍须在教育部的推免系统中完成填报志愿等相关操作，否则申请无效”（\nref{S0173} 中国矿业大学；\nref{S0170} 沈阳药科大学同义）。天津财经大学也说“已在我校预报名系统中报名的学生，后期须获得本校推免资格，并在研招网推免系统进行正式报名后，且接受我校待录取，方可视为录取成功”（\nref{S0416}）。所以正确理解是：校内或预推免系统决定“能不能进复试”，全国推免服务系统决定“录取是否成立”。
\nref{S0173}
\nref{S0170}
\nref{S0416}
\end{qa}

\iin{中国矿业大学}
\ith{两个系统}
\ist{九推与系统填报}
\begin{qa}{E0036}{学校自己的预推免系统和推免服务系统是什么关系？谁说了算？}
全国推免服务系统是唯一权威，校级系统只是前置筛选。研招网对系统的定位是“全国统一的推免工作信息备案公开平台和推免生网上报考录取系统”，资格审核确认、报考、录取以及备案公开均须通过该系统进行，且“最终推免生名单、录取名单和推免生学籍注册均以推免服务系统备案信息为准”（\nref{S0464}）。天津大学的说法更直白：“9 月 21 日教育部‘推免申请系统’开通之前，所有申请我校的推免生，均须通过我校推免系统填报预报名”（\nref{S0520}）。结论：校级系统决定你能不能拿到考核与复试机会，全国系统决定录取是否生效（\nref{S0418}）。
\nref{S0464}
\nref{S0520}
\nref{S0418}
\end{qa}

\iin{教育部学生服务与素质发展中心（学信网/研招网）}
\ith{系统关系}
\ist{九推与系统填报}
\begin{qa}{E0037}{夏令营拿到的优秀营员，九推还要不要在全国系统里再走一遍流程？}
要走，而且是必须走。研招网明确“所有推免生均通过‘推免服务系统’填写报考志愿、接收并确认招生单位的复试及待录取通知”（\nref{S0464}），且最终推免生名单、录取名单和学籍注册均以该系统备案信息为准。实际操作上只是简化：四川大学要求获得推免资格且符合条件的考生“须在国家规定时间内登录推免服务系统完成报名、填报志愿”（\nref{S0245}）；武大计算机学院要求通过预报名者“须于全国推免服务系统开放后的 5 小时内填报我院”（\nref{S0119}）。另外提醒一句：并非所有夏令营都评选优秀营员，华东师大信息与电子工程学院就明确“本次夏令营不评选优秀营员”（\nref{S0217}）。
\nref{S0464}
\nref{S0245}
\nref{S0119}
\nref{S0217}
\end{qa}

\iin{教育部学生服务与素质发展中心（学信网/研招网）}
\ith{优营衔接}
\ist{九推与系统填报}
\begin{qa}{E0038}{学校说“优秀营员九推直接填系统就行”，这是真的吗？还要复试吗？}
大体为真，但流程里仍有必须点的两步。以上海财经大学为例：“已通过我校推荐免试预报名选拔获得我校拟录取资格的推免生，请于 9 月 21 日 14:00 前……填报已被我校拟录取的专业志愿，逾期未填报专业志愿将视为放弃。按照教育部规定，我校将于 9 月 22 日上午向有关同学发放复试通知，待推免生接受我校复试通知后，我校将于 9 月 24 日上午向有关推免生发放待录取通知”（\nref{S0183}）。也就是说复试多为形式性审查，但“填报志愿—接受复试通知—接受待录取”三步一个都不能漏，漏掉任何一步都可能被视为放弃。
\nref{S0183}
\end{qa}

\iin{上海财经大学}
\ith{优营衔接}
\ist{九推与系统填报}
\section{技术故障}
\begin{qa}{E0005}{照片上传一直提示格式或尺寸不合格，到底要求什么规格？手机裁的图能用吗？}
可参考上海财经大学公布的推免服务系统照片要求：照片大小应介于 1KB 至 100KB；高度应在 100 至 720 像素之间；宽度应在 90 至 480 像素之间，且小于高度；格式为 jpg（\nref{S0182}）。手机直接裁的图常见问题是：文件大于 100KB、宽高比不满足“宽小于高”、或导出成 png、heic 格式。处理办法是导出为 jpg 后按上述区间压缩并裁成竖向比例，再对照系统页面提示复核（\nref{S0182} 同时提醒“注意查看系统页面的相关说明”）。若招生单位另行给出照片规格，以该校当次通知为准。
\nref{S0182}
\end{qa}

\iin{上海财经大学}
\ith{照片上传}
\ist{九推与系统填报}
\begin{qa}{E0006}{学籍校验一直不通过，提示与学信网信息不一致，该怎么办？会不会耽误填志愿？}
按官方指引找源头改，不要只找系统客服。研招网在推免系统报考注意事项中写明：“推免生如对系统自动填入的基本学籍（学历）信息有疑问，请尽快向就读高校的学籍管理部门咨询”（\nref{S0389}、\nref{S0444}）。系统里自动填入的学籍信息来自学信网学籍库，因此不一致多半是学籍库中的姓名、证件号、学号、预计毕业时间等字段与你的实际不符，必须在学校侧更正。时间上不必过度恐慌：个人资料填报自 9 月 18 日开放（\nref{S0464}），填报志愿 9 月 21 日 9:00 才开始，只要在填志愿前处理完即可；实在来不及，先把能填的志愿填上，同时推进学籍更正。
\nref{S0389}
\nref{S0444}
\nref{S0464}
\end{qa}

\iin{www.sceea.cn}
\ith{学籍校验}
\ist{九推与系统填报}
\begin{qa}{E0032}{系统高峰期一直转圈、报错 502，是换时间点还是换浏览器？}
两者都值得做，但官方没有“高峰时段”的说明，能确认的是浏览器建议。北京大学多个院系要求“请使用 Chrome 或 Firefox 浏览器”（\nref{S0114}、\nref{S0299}）；四川大学预推免平台提示“如使用过程出现异常，推荐使用 Chrome 或者火狐浏览器重试”（\nref{S0474}）；深圳大学“建议使用 Google Chrome 浏览器或 360 极速浏览器”（\nref{S0241}）；华中科技大学“推荐在电脑上使用 Chrome 或 360 浏览器极速模式”（\nref{S0198}）。至于转圈、报错 502 这类现象，未见研招网公开说明，属经验范畴。实务上先换 Chrome 或 Edge 的无痕模式并关闭拦截类插件，再错开 9:00、14:00 等开放时点操作。
\nref{S0114}
\nref{S0299}
\nref{S0474}
\nref{S0241}
\nref{S0198}
\end{qa}

\iin{北京大学}
\ith{系统故障}
\ist{九推与系统填报}
\begin{qa}{E0033}{只用手机浏览器填报行不行？是不是必须用电脑？}
未见官方“必须用电脑”的规定，但院校建议清一色指向桌面浏览器。北京大学多个院系要求“请使用 Chrome 或 Firefox 浏览器”（\nref{S0114}、\nref{S0299}）；华中科技大学“推荐在电脑上使用 Chrome 或 360 浏览器极速模式”（\nref{S0198}）；中国农业大学博士网报“建议使用 Windows 系统 360 浏览器，极速模式”（\nref{S0160}）；四川大学推荐“Chrome 或者火狐浏览器”（\nref{S0474}）。这些提示都默认桌面环境。经验上手机端在高峰时段更容易卡在提交与支付环节。稳妥做法是：志愿提交、支付与待录取确认用电脑完成，手机只用于查看通知和接收短信。
\nref{S0114}
\nref{S0299}
\nref{S0198}
\nref{S0160}
\nref{S0474}
\end{qa}

\iin{北京大学}
\ith{设备要求}
\ist{九推与系统填报}
\begin{qa}{E0059}{系统崩溃导致我错过了确认时限，这种情况能申诉吗？}
可以申诉，但能否挽回取决于招生单位，且官方并未给出“系统故障免责”的规则。研招网给出的申诉方向是：对备案公开的推免招生信息有疑问，“可直接与相关推荐高校、招生单位联系咨询，或向推荐高校、招生单位所在省级教育行政部门、省级教育招生考试机构申诉”（\nref{S0464}）；系统内每个招生单位页面也公示了申诉渠道与联系方式（\nref{S0465}）。对统考调剂，教育部明确规定招生单位发布需考生确认的通知时“应充分考虑考生学习、工作、休息时间等因素作出合理安排，给考生预留充裕的确认时间”，不得简单以“逾期不接受视为自行放弃”对待（\nref{S0049}），这条精神可用于申诉。第一时间要留存系统报错截图与操作时间，并记录网址、浏览器版本、网络环境、通知类型和截止时刻；在仍有时间时换稳定网络与桌面浏览器重试，但不要连续刷新导致状态不明。随后用官方电话和邮件同时联系招生单位，邮件主题写明姓名、报考院系专业、通知类型与时间，附截图并请求确认是否可补操作；必要时再向推荐高校及省级教育招生考试机构反映。保存发送记录和回执。系统故障不自动延长个人时限，申诉渠道也不保证恢复资格。
\nref{S0464}
\nref{S0465}
\nref{S0049}
\end{qa}

\iin{教育部学生服务与素质发展中心（学信网/研招网）}
\ith{系统故障申诉}
\ist{九推与系统填报}
\section{报名费}
\begin{qa}{E0028}{推免报名费是多少？是每填一个学校交一次，还是只交一次？}
没有全国统一金额，标准由推荐高校所在省级教育招生考试机构确定。研招网《2027推免生网上支付报名费注意事项》原文：“推免服务系统提供代收推免报名费服务，报名费标准和规则均以推荐高校所在省级教育招生考试机构规定为准。推免生对报名费如有疑问，请咨询推荐高校所在省级教育招生考试机构”（\nref{S0451}）；河北大学的通知也写“报名费标准以省级教育招生考试机构规定的为准”（\nref{S0604}）。缴费在流程上位于“填报个人资料”之后、填报志愿之前（\nref{S0389}、\nref{S0444}），因此全流程只缴一次，不是每填一个志愿各缴一次；系统关闭注册、支付、报名功能的时间为 10 月 18 日（\nref{S0604}）。
\nref{S0451}
\nref{S0604}
\nref{S0389}
\nref{S0444}
\end{qa}

\iin{教育部学生服务与素质发展中心（学信网/研招网）}
\ith{报名费标准}
\ist{九推与系统填报}

\begin{qa}{N0001}{2027 届推免生的录取信息什么时候公开？系统暂时没有日期怎么办？}
2027 届系统说明目前给出了注册、填报志愿、确认复试和待录取的时间，但对“录取信息公开”只写明请关注网站，尚未列出该届统一公示起止日期。不要把某一院校的拟录取公示日当成全国最终备案日。先看目标单位官网和推免服务系统公告；临近入学仍查不到本人备案状态时，分别联系招生单位与本校教务部门核对。网页日期会更新，操作前应重新查看研招网当页。\nref{S0761}
\qscope{全国推荐免试攻读研究生信息公开管理服务系统2027届说明，页面信息核对于2026年10月12日；公示日以之后更新的官方通知为准。}

查询时把“院校拟录取公示”与“全国系统录取信息公开”分开记日期：前者由学校安排，不能据此推断后者已完成。可以每周查看研招网推免系统首页/信息公开栏目和目标单位研究生院公告，记录页面更新时间；若公示期已结束而本人系统状态仍有疑问，向招生单位研招办提供考生号、专业和系统状态截图，问清是否已完成备案以及何时能查询。2027届统一公开日期尚未列出时，不要转用2026届日期作承诺。
\end{qa}

\begin{qa}{N0002}{研招网展示的同专业录取公示，能用来推算完整名额和上岸率吗？}
不能直接这样推算。系统说明限定为“本人报考并参加复试的招生单位同专业、同层次、同类别”且只在规定时限内查看；它不是所有学校、所有批次和所有培养类别的汇总统计。名单还可能受单位分批公示、专项计划单列和专业名称口径影响。可把它用于核对公开结果，估算时还要查当年招生目录、学院接收办法和最终公示名单，并注明统计范围与截止日期。\nref{S0761}
\qscope{全国推免服务系统2027届说明；名单展示范围依官方页面规定。}

做任何比例估算前先固定分母：招生目录计划数、学院实际接收数、系统当前可见名单不是同一个口径。把硕士/直博、普通/专项、学院方向以及名单发布日期分别列开；若计划调整或分批公示，应优先使用最终公示和学校当年接收办法，并写明你只统计了哪些公开样本。样本不足时只把名单用于了解录取结果，不计算“上岸率”，也不要把可见人数反推成剩余名额。
\end{qa}

\begin{qa}{N0003}{同专业的公示名单里为什么看不到其他学校或其他方向的录取者？}
系统公开范围不是全网搜索：官方说明只允许在规定时限内查看本人曾报考并参加复试的招生单位中，同专业、同层次、同类别推免生的录取公示。未参加该单位复试、不同层次或类别的考生，不属于这一查询范围；方向名称也可能合并展示或另行公示。需要比较不同方向时，应回到各院系官网的拟录取名单和专业目录，不能把系统未显示理解为没有招生。\nref{S0761}
\qscope{全国推免服务系统2027届说明，适用范围以系统页面和招生单位实际公示为准。}

核对时可先检查四个限定项：招生单位是否相同、专业是否相同、硕士/直博层次是否相同、普通/专项类别是否相同；任一不同都可能不在同一查询列表。学院官网还可能把研究方向合并展示，或分批发布更正名单。若用于个人确认，核对自己的姓名、专业和类别即可；若用于统计，不要把看不见的考生当成未录取，也不要截图传播他人的身份信息。
\end{qa}

\begin{qa}{N0004}{学校官网已经公示拟录取名单，为什么推免系统里还没有最终备案信息？}
官网拟录取公示与系统备案属于不同环节。招生单位先按本单位办法复试、审定并公示名单，之后按规定在推免服务系统完成备案；研招网明确最终录取名单和学籍注册以系统备案信息为准。看到官网名单后，继续按系统通知完成待录取确认，并保存公示页面和系统状态截图；若两边长期不一致，联系招生单位研招办核实，不要自行认定哪一条已经生效。\nref{S0761}
\qscope{全国推免服务系统2027届说明；各校具体审核和备案时点不同。}
\end{qa}

\begin{qa}{N0005}{推免系统显示录取信息后，还需要核对哪些字段？}
核对招生单位、院系、专业、学习层次、培养类型和个人身份信息是否与已确认的志愿一致；直博生还要确认方向或导师等单位要求的字段。随后下载或截图保存系统页面、单位公示和通知时间，发现错项立即通过官方渠道联系招生单位及本校。系统最终备案信息关系到学籍注册，不能只凭口头承诺或聊天记录。\nref{S0761}
\qscope{全国推免服务系统2027届说明；具体字段按系统当年页面和招生单位通知核验。}

建议把系统记录与正式录取通知、学校公示逐项对照，至少检查姓名及证件信息、招生单位/院系、专业名称与代码、硕士或直博、全日制/非全日制、定向/非定向及专项类别；存在研究方向或导师字段时也一并确认。随后保存带页面来源与时间的本人记录，遮盖证件号后再用于咨询。发现差异时把字段名、系统显示值、应有值及支持文件一次发给研招办，询问是否为展示口径差异还是备案错误，并保存书面处理结论。
\end{qa}
\begin{qa}{E0029}{报名费交了但面试没去，能退吗？}
未见研招网与学信网关于推免报名费退费的官方说明，无法给出确定答复。能找到的是两点：一是报名费标准和规则均以推荐高校所在省级教育招生考试机构规定为准，有疑问请咨询该省级机构（\nref{S0451}）；二是院校侧的同类费用普遍不退——北航仪器科学与光电工程学院明确“因各种原因未参加复试者，已支付的复试费不予退还”（\nref{S0421}），北京科技大学冶金与生态工程学院也是“如因个人原因未按时参加复试的，已缴复试费不予退还”（\nref{S0270}），清华大学经济管理学院写明申请材料和已缴纳的复试费不予退还（\nref{S0254}）。建议缴费前就先向省级机构问清退费口径。
\nref{S0451}
\nref{S0421}
\nref{S0270}
\nref{S0254}
\end{qa}

\iin{教育部学生服务与素质发展中心（学信网/研招网）}
\ith{报名费退费}
\ist{九推与系统填报}

\section{2027届窗口末段操作（2026年10月16日核对）}
\begin{qa}{E0491}{今天已经拿到一个待录取，系统里应该先核对什么再确认？}
先别只看学校简称或聊天里的“拟录取”。打开通知详情逐项核对招生单位、院系、专业代码与名称、硕士/直博类别、学习方式以及回复截止时刻；把这些字段与本人志愿、招生单位正式通知对照。确认期限短于全国系统关闭时间时，仍按通知的个人截止时间处理。存在专业方向或培养类别不一致，先立即联系招生单位研招办确认，并保留通知页面；不要在信息未核实时点击接受。最终录取与学籍以推免服务系统备案为准。\nref{S0761}\nref{S0761}
\qscope{全国推免服务系统2027届说明；2026年10月16日核对，个案时限以系统和招生单位通知为准。}
\end{qa}

\begin{qa}{E0492}{2027届10月18日填报截止后，10月19日还能新填志愿吗？}
不能把10月20日系统关闭理解成还能补报新志愿。研招网说明把2027届填报窗口列为2026年9月21日至10月18日，系统整体关闭日为10月20日；两者处理的事项不同。10月18日后不要假设可新建或修改志愿，已提交志愿衍生的复试、待录取通知则按系统页面和招生单位通知及时处理。若页面功能与公告理解不一致，以当年系统实际提示并向研招网/招生单位核实。\nref{S0761}\nref{S0761}
\qscope{全国推免服务系统2027届说明；2026年10月16日核对，个案时限以系统和招生单位通知为准。}
\end{qa}

\begin{qa}{E0493}{只收到电话或微信口头“拟录取”，系统还没通知，算录取了吗？}
不算已经完成系统确认。口头消息可以作为及时准备的信号，但不是系统内的复试或待录取通知，也不能代替你在规定时限内作出操作。请通过招生单位官网公布的电话或邮箱核实是否已发送、是否需要补材料以及最晚确认时间；同时持续查看推免系统站内通知和绑定联系方式。不要因为口头承诺停止处理其他尚未接受的流程，接受待录取后系统规则会限制后续志愿与通知。\nref{S0761}\nref{S0761}
\qscope{全国推免服务系统2027届说明；2026年10月16日核对，个案时限以系统和招生单位通知为准。}
\end{qa}

\begin{qa}{E0494}{复试通知和待录取通知同时到达，应该先处理哪一个？}
逐条看清通知类型、截止时间和后果，不按短信到达顺序猜。复试通知对应是否参加该单位复试，待录取通知则进入只能接受一个的关键选择阶段；系统说明要求按时接收并确认招生单位通知，接受待录取后不能再填报其他志愿或接受其他单位相关通知。先记录每条通知的单位、专业、类型和截止时刻；如果待录取通知期限更早，且你尚未完成比较，立即联系该单位确认能否延时，同时联系其他单位核实考核结果，不能把未回复视为自动延期。\nref{S0761}\nref{S0761}
\qscope{全国推免服务系统2027届说明；2026年10月16日核对，个案时限以系统和招生单位通知为准。}
\end{qa}

\begin{qa}{E0495}{志愿锁定期间收到更想去的学校通知，能不能直接换志愿？}
一般不能直接改掉已提交且仍处于锁定期的志愿。先在系统查看该志愿的提交时间和当前状态，再确认是否已收到复试通知、单位是否明确拒绝，以及系统页面是否已开放修改。不要通过重复注册或换浏览器尝试绕过锁定，也不要为了腾出名额轻率拒绝已收到的通知；拒绝可能不可撤回。若确有特殊情况，联系招生单位确认本单位处理方式，并以系统实际状态为准。时间窗口和志愿规则按2027届研招网说明执行。\nref{S0761}\nref{S0761}
\qscope{全国推免服务系统2027届说明；2026年10月16日核对，个案时限以系统和招生单位通知为准。}
\end{qa}

\begin{qa}{E0496}{系统里的复试确认截止时间和学院邮件不一样，怎么避免错过？}
把两处截止时间分别记录，先按更早的时刻行动，并通过学院官网公布的官方联系方式确认是否为时区、批次或录入差异。保存邮件原文、系统通知详情和咨询答复；若对方确认以其中一个为准，请尽量取得书面回复。全国系统提供统一报考录取流程，但各招生单位会发布本单位安排，不能把系统关闭日当成每条通知的个人答复期限。确认后仍回到系统完成必要操作并检查状态是否更新。\nref{S0761}\nref{S0761}
\qscope{全国推免服务系统2027届说明；2026年10月16日核对，个案时限以系统和招生单位通知为准。}
\end{qa}

\begin{qa}{E0497}{待录取通知点了“拒绝”或错过时限，还有撤回机会吗？}
不要预设可以恢复。系统说明对通知确认和志愿解锁有明确状态规则，但未承诺误拒或超时后可撤销；是否还能补救取决于招生单位是否仍有名额、系统是否允许重新发送。发现误操作后，立刻截图记录通知、操作时间和当前状态，通过招生单位研招办官方联系方式说明情况，并询问能否重新发通知，同时联系本校推免管理部门。不要等到次日，也不要把电话沟通当成系统状态已恢复。\nref{S0761}\nref{S0761}
\qscope{全国推免服务系统2027届说明；2026年10月16日核对，个案时限以系统和招生单位通知为准。}
\end{qa}

\begin{qa}{E0498}{已经接受一所学校待录取，其他学校还在等复试结果，能继续等吗？}
接受待录取不是“先占位、再比较”的可逆动作。研招网说明明确，接受招生单位待录取后不能再填报其他志愿或接受其他招生单位的复试、待录取通知，其他单位也不能再向你发送相关通知。接受前把培养层次、专业方向、导师沟通、地点、学制与费用等已核实信息放在同一张比较表，并确认自己愿意承担放弃其他机会的后果；尚有关键问题未确认时，先联系发通知单位询问期限，不要用“先接受再说”的方式占位。\nref{S0761}\nref{S0761}
\qscope{全国推免服务系统2027届说明；2026年10月16日核对，个案时限以系统和招生单位通知为准。}
\end{qa}

\begin{qa}{E0499}{系统显示“已接受”后还要不要截图和下载通知？}
要。至少保存通知详情、确认状态、操作时间、招生单位公示页面和相关邮件；文件名写明日期、学校、院系、通知类型，另记下官方咨询渠道及答复。截图应包含页面标题、状态与时间，避免只截到无法识别来源的一小块；重要文件同步到两个安全位置。留存是为了解决后续状态核对，不代表截图能替代系统备案。最终名单与学籍注册依据仍是推免服务系统备案信息。\nref{S0761}\nref{S0761}
\qscope{全国推免服务系统2027届说明；2026年10月16日核对，个案时限以系统和招生单位通知为准。}
\end{qa}

\begin{qa}{E0500}{10月18日至20日只需要等系统自动处理吗？}
不应被动等待。10月18日是2027届志愿填报截止，10月20日是系统关闭节点；等待期间仍要主动查看已提交志愿的复试和待录取通知、处理个人答复期限、联系尚未反馈的招生单位，并确认系统状态与官网名单是否一致。提前检查绑定手机号、常用邮箱和登录方式；每次操作后刷新确认状态并保存凭据。全国系统没有替你自动接受通知的机制，临近关闭时应以页面显示和单位正式通知为准。\nref{S0761}\nref{S0761}
\qscope{全国推免服务系统2027届说明；2026年10月16日核对，个案时限以系统和招生单位通知为准。}
\end{qa}
